\section{Introduction}
\label{sec:intro}

A superoptimizer searches for the best program computing a given function. For most of the field's history the claim it made was modest: \emph{here is a better program than the compiler emitted}, and that claim is cheap to check---replay the program. Recent systems make a stronger claim. HieraSynth~\cite{hierasynth} states that its completeness, \enquote{combined with an accurate cost model, proves that the generated program is optimal, assuring experts that further manual tuning would be futile.} Optimality is a non-existence claim over an entire search space, and it is established the same way in every system from Massalin's 1987 superoptimizer~\cite{massalin} to the present: the space is exhausted, a solver answers \emph{unsat} at each leaf, and the refutation is discarded. The proof, when a reader asks for one, is the solver returning \emph{unsat} with the timeout switched off.

This is not a criticism of those systems. It is a description of a gap between two communities that have not met. The SAT community has required machine-checkable proofs of every \emph{unsat} in its main competition track since 2013, disqualifies solvers whose certificates fail to check, and has formally verified checkers for them~\cite{cakelpr}. The compiler community certifies \emph{correctness}---that a rewrite preserves semantics---with translation validation~\cite{alive2} and verified compilers~\cite{compcert}, and deliberately excludes \emph{cost} from the proof: one recent verified scheduler \enquote{abstracts away the outside scheduling heuristic as a universal parameter so it is flexible to modify without touching any correctness proof}~\cite{yang2024}. Correctness is a universal statement discharged once per rewrite. Minimality is a non-existence statement over a search space, and non-existence is exactly what a refutation proof is for. Nobody had joined the two.

We describe a reasoning stack, axeyum, built on the principle that search is untrusted and checking is small, and we use it to produce what we believe are the first per-instance machine-checkable minimality certificates for instruction sequences on a current ISA. The instance is deliberately elementary---reversing the thirty-two bytes of an AVX2 register---because the point is the artifact, not the theorem: every SIMD programmer knows the two-instruction answer, LLVM emits it, and nobody could previously state that one instruction does not suffice as anything but folklore, because nobody had written down the instruction set the claim quantifies over.

\subsection{Contributions}

\begin{itemize}
\item \textbf{An encoding that makes minimality pure SAT.} Lanes are modelled as distinct symbolic provenance tags rather than bytes, so a permutation-only instruction set needs no bit-vector arithmetic; each instruction is a finite-domain function on tag vectors, and a proof-producing CDCL core emits DRAT natively (\cref{sec:method}).
\item \textbf{Three certified theorems with declared scope.} Byte reversal has minimum length two in a unary five-family AVX2 language and in a fourteen-family SSA language, and minimum cost four under a named Haswell dependent-latency profile; certificates of 1--13\,MB, re-checked by an independent backward checker, with truncation controls (\cref{sec:results}).
\item \textbf{A vacuity finding.} One of the three certificates is accepted by the checker with its first half deleted, because its formula is refutable by unit propagation alone. The checker is sound; the certificate carries no information. We give the general statement, a one-line test, and argue that every small certificate should be subjected to it (\cref{sec:vacuous}).
\item \textbf{Reproductions of the uncertified searches behind a ratified ISA.} We re-derive, exactly, the 256-function table on which RISC-V's scalar-crypto extension answered a design question, and the permutation-reachability table whose source artifact is now a dead link; and we identify a bare lower bound asserted in a specification with no proof (\cref{sec:isa}).
\end{itemize}

We are careful about what is new. Bansal and Aiken enumerated optimal three-instruction x86 peepholes in 2006~\cite{bansalaiken}; Franchetti and Püschel proved stride-permutation sequences minimal on SSE2 and Cell in 2008 by matching a search's upper bound to Floyd's block-transfer lower bound~\cite{franchetti2008}. We are not first to claim instruction-sequence optimality. We are first, to our knowledge, to make the claim re-checkable by a third party without re-running or trusting the tool, for general permutations on a current ISA, with a certificate per instance.
